% Figure: the conduction shape factor for a cylinder of radius r
% buried at depth z below an isothermal ground surface. Isotherms of
% the harmonic potential are eccentric circles (bipolar coordinates);
% the surface is flat, the pipe is warm, and heat spreads outward.
% The shape factor S packs this whole field into a single conductance.
\begin{figure}[htbp]
\centering
\begin{tikzpicture}[>=Stealth, font=\small]
% Ground surface (isothermal, T_s), a horizontal line.
\draw[engblack, very thick] (-4.2,0) -- (4.2,0);
\node[anchor=south west, engblack] at (-4.2,0) {surface at $T_{s}$};
% Hatch above surface to mark the air / boundary side.
\foreach \x in {-4.0,-3.5,...,4.0}{
  \draw[engblack!45, thin] (\x,0) -- (\x+0.28,0.28);
}
% The buried pipe: a filled circle at depth z.
\fill[engred!22] (0,-2.4) circle (0.55);
\draw[engred, very thick] (0,-2.4) circle (0.55);
\node[engred] at (0,-2.4) {$T_{p}$};
% Depth annotation z from surface to pipe centre.
\draw[<->, engblack] (2.4,0) -- (2.4,-2.4);
\node[anchor=west, engblack] at (2.4,-1.2) {depth $z$};
% Radius annotation.
\draw[<->, engblue] (0,-2.4) -- (0.55,-2.4);
\node[anchor=north, engblue] at (0.30,-2.55) {$r$};
% Isotherms: eccentric circles crowding the pipe, widening toward the
% surface. Drawn as ellipse-like arcs at increasing radius and rising
% centre to suggest the bipolar family.
\draw[engblue!70, thick]   (0,-2.4) ellipse (0.95 and 0.85);
\draw[engblue!55, thick]   (0,-2.15) ellipse (1.55 and 1.30);
\draw[engblue!40, thick]   (0,-1.75) ellipse (2.35 and 1.70);
\draw[engblue!30, thick]   (0,-1.30) ellipse (3.30 and 1.10);
\node[anchor=west, engblue!60] at (1.7,-2.35) {isotherms};
\end{tikzpicture}
\caption[Conduction shape factor for a buried cylinder]{Steady heat
conduction from an isothermal cylinder of radius $r$ buried at depth
$z$ below an isothermal ground surface, a two-dimensional Laplace
problem on the semi-infinite domain. The isotherms (blue) are the
level sets of the harmonic potential; they crowd against the warm pipe
and spread as eccentric circles toward the cool surface, the picture
that the bipolar-coordinate solution makes exact. The entire field
collapses into a single number, the conduction shape factor
$S = 2\pi/\operatorname{arccosh}(z/r)$, so that the total heat rate is
$q = S k (T_{p} - T_{s})$: a distributed Laplace field reduced to one
lumped conductance $Sk$. The engineer sizing a buried steam line or a
ground-loop heat exchanger reads the loss from $S$ alone, with no
field solve per case.}
\label{fig:vol02:ch12:shape-factor}
\end{figure}
